\chapter{Implementation}
\label{ch:impl}

This chapter describes the harness as built on 9 October 2026 (milestones M0--M2 complete, M3 steps 1--5 complete).
The description follows \code{docs/architecture.md}, which is kept up to date with the code.

\section{Stack and code organisation}

The harness is written in Python (3.11 or newer) and managed with \code{uv}. Model calls go through LiteLLM
\cite{litellm}; tool arguments and answers are Pydantic models; graph algorithms come from networkx \cite{networkx};
the MCP server uses the official \code{mcp} Python SDK \cite{mcp}; model-written code runs in Docker. Models currently
run locally through Ollama, so there are no API keys in the loop yet; any secrets are read from a git-ignored
\code{.env} file.

\begin{table}[ht]
  \centering
  \small
  \caption{Modules of the harness and the evaluation code.}
  \label{tab:modules}
  % source: docs/architecture.md
  \begin{tabularx}{\textwidth}{>{\raggedright\arraybackslash}p{4.6cm}>{\raggedright\arraybackslash}X}
    \toprule
    File & Job \\
    \midrule
    \code{harness/loop.py} & The agent loop: \code{run_agent}, \code{AgentConfig}, \code{RunResult} \\
    \code{harness/models.py} & One LiteLLM call: tokens, latency, retries, context window \\
    \code{harness/tools/graph_tools.py} & Graph tools, argument models, generated schemas, \code{run_tool} \\
    \code{harness/tools/handles.py} & Result handles and \code{read_result} \\
    \code{harness/tools/mcp_server.py} & The same tools over MCP (stdio) \\
    \code{harness/sandbox/} & \code{run_python}: tool schema, Docker backend, in-container runner, image \\
    \code{harness/compaction.py} & Shortening older messages near the context window \\
    \code{harness/brief.py} & Short versions of values repeated back to the model \\
    \code{harness/answers.py} & Answer types, \code{submit_answer}, \code{cannot_answer}, shape checks \\
    \code{harness/verifiers.py} & Evidence-based checks of final answers \\
    \code{harness/tasks.py} & Task registry: template, answer type, verifier \\
    \code{harness/parsing.py} & Spotting and rescuing tool calls written as text \\
    \code{harness/prompts.py} & All model-facing text and \code{PROMPT_VERSION} \\
    \code{harness/trace.py} & Append-only JSONL trace \\
    \code{eval/tasks.py} & Dataset loading, generated questions, reference answers, grading \\
    \code{eval/runner.py} & Runs the agent over questions, writes \code{results.jsonl} and traces \\
    \code{eval/analysis/report.py} & Accuracy, confidence intervals, tokens, time, code use \\
    \code{eval/analysis/process.py} & Tool precision, recall, F1, parameter match \\
    \code{eval/analysis/trajectory.py} & Classification of how a run reached its answer \\
    \bottomrule
  \end{tabularx}
\end{table}

\section{One question, end to end}
\label{sec:one-question}

The evaluation runner loads a question and its graph, looks up the task in the harness's registry (template, answer
type, verifier) and calls \code{run_agent}. The loop then proceeds as follows.

\begin{enumerate}
  \item The conversation starts with the system prompt and the question. The tools offered are the graph tools, plus
        \code{submit_answer} and \code{cannot_answer}.
  \item Up to the step limit, the model is called with the conversation and the tools.
  \item The model's hidden thinking is removed from its reply before the reply is appended to the history, and logged
        separately.
  \item If the reply contains a tool call written as plain text, the harness either rescues it (parses and runs it,
        lenient mode) or answers with a format error (strict mode). Rescues are counted and logged either way.
  \item If the same reply (text and tool calls, ignoring call identifiers) arrives three times in a row, the run stops
        with status \code{loop_detected}.
  \item Each tool call is handled in turn: arguments that are not a JSON object come back as an error; a graph tool
        is run through \code{run_tool}, which validates the arguments; large values become handles;
        \code{read_result} pages through a handle; \code{submit_answer} resolves any handles, checks the answer's
        shape and then runs the verifier, which either accepts the answer or sends an error back; \code{cannot_answer}
        ends the run.
  \item A reply with no tool call gets a nudge: a format error if it looks like a hand-written tool call, an
        empty-reply message if it has no content at all, and a general nudge otherwise.
  \item Every step is logged to the trace.
\end{enumerate}

The loop returns a \code{RunResult} with the answer, the status (\code{submitted}, \code{cannot_answer},
\code{max_steps}, \code{loop_detected} or \code{context_full}), the evidence, and counters for steps, tokens, model
latency, rescues, rejections and compactions. The runner grades the answer against the reference with
\code{eval/tasks.py}, independently of the verifier, and writes one row to \code{results.jsonl}.

\section{The agent loop}

\subsection{Configuration}

An \code{AgentConfig} holds the settings that stay fixed across the questions of one experiment: the model, the step
limit (default 10), the repeat limit for loop detection (default 3), lenient or strict handling of text tool calls,
which graph tools are offered, which tool functions are available inside code, the handle threshold (default 200
numbers), whether and how \code{run_python} is offered (\code{None}, \code{"tools"} or \code{"networkx"}), whether
the code strategy hint is added to the system prompt, and the compaction threshold (default 75\% of the context
window).

\subsection{Failure handling found by poking the loop}

Several mechanisms in the loop came from early failures of a naive version, observed with qwen3:8b on small graphs
(learning log, 2026-10-04):
\begin{itemize}
  \item Asked for the degree of a node that does not exist, the model read the tool error and then submitted 0,
        because \code{submit_answer} only accepted an integer. The schema shaped the behaviour; the fix was an honest
        exit, \code{cannot_answer(reason)}, logged as its own status rather than as a wrong answer.
  \item Asked for the largest clique with no clique tool, the model first said that the tools were insufficient,
        was nudged twice with a generic ``use the tools or submit'' message, and then guessed. The harness itself had
        caused a hallucination; the nudge now also offers \code{cannot_answer}.
  \item The model wrote \texttt{submit\_answer \{"answer": 0\} </tool\_call>} as plain text nine times until the step
        limit. Feedback now distinguishes a hand-written tool call (\code{FORMAT_ERROR}) from other text
        (\code{NUDGE}), and loop detection ends a run after three identical replies.
\end{itemize}
After these fixes, the clique question ended in \code{cannot_answer} at step~1, and the missing-node question in
\code{cannot_answer} at step~2. The degree question about a missing node gave \code{cannot_answer} in two of four
runs and 0 in the other two: improved, not solved, and a first reminder that one run is not a result.

\subsection{Rescuing text tool calls}

Small models sometimes write a tool call into their text instead of using the function-calling interface. The
harness detects this (tool-call tags, or a tool name followed by a JSON object) and, in lenient mode, parses and runs
the call. Each rescue is logged as a separate event, so that every experiment reports two accuracies: lenient (rescued
answers count) and strict (an answer that needed a rescue counts as wrong). Prose that only mentions a tool name is not
flagged. In the results so far this is the main source of the gap between the two accuracies
(Chapter~\ref{ch:results}).

\subsection{Thinking is not re-sent}

Reasoning models such as qwen3 return their hidden thinking in a separate field of the reply. In an early version the
loop kept that field in the history, so the thinking was sent back and re-read on every later step. Removing it before
the reply joins the history reduced the second-step prompt of one question from 1,826 to 534 tokens. Results before
this fix have inflated prompt tokens on steps after the first; accuracy and output tokens are unaffected.

\section{Tool layer}

\subsection{Tools as validated functions}

Each tool is a plain function \code{tool(graph, **args) -> dict} registered once, together with a Pydantic model for
its arguments and the description the model reads. The JSON schemas shown to the model are generated from the
argument models, and \code{run_tool} validates incoming arguments with the same models, so the schema the model reads
and the validation the harness runs cannot drift apart. When this was introduced, the generated schemas were checked to
be byte-identical to the earlier hand-written ones, so earlier results remained comparable.

Validation produces messages the model can act on, for example ``Bad arguments for shortest\_path: target: Field
required'' or ``node: a node id must be an integer, not true/false''. Two details matter in practice: the string
\code{"3"} is accepted as node 3 (before, it was looked up as a string and reported as missing), and \code{true} is
rejected as a node identifier (Pydantic's lax integer type would have turned it into 1, so node identifiers use a
custom type that refuses booleans). The same strictness applies to answers: in Python a boolean is an integer, so a
number answer rejects booleans and a yes/no answer rejects 0 and 1.

\subsection{The tool set}

There are 13 default tools: \code{get_neighbors}, \code{degree}, \code{has_edge}, \code{graph_info},
\code{shortest_path}, \code{connected_components}, \code{has_cycle}, \code{minimum_spanning_tree}, \code{has_path},
\code{is_bipartite}, \code{topological_sort}, \code{distances_from} and \code{neighborhood}. Five further tools are
opt-in, offered only when requested: \code{articulation_points}, \code{bridges}, \code{k_core}, \code{triangles} and
\code{greedy_coloring}. They came out of the GDS Agent review and are opt-in because every offered tool changes the
prompt, and because \code{triangles} would answer one of the combine tasks in a single call. Appendix~\ref{app:tools}
lists all tools with their arguments and descriptions. New tools are appended, so the text of earlier tools stays the
same.

Where possible, tools return the evidence for their answer, not just the answer: \code{has_cycle} returns one cycle,
\code{has_path} returns a shortest path, \code{is_bipartite} returns the two sides or an odd cycle (so both answers
carry a proof), and \code{topological_sort} returns an order or a directed cycle. The odd cycle comes from the
breadth-first coloring itself: when an edge joins two nodes of the same color, walking both up the BFS tree to their
meeting point gives an odd cycle. When networkx signals ``nothing found'' by raising an exception (no path, no cycle),
the tool turns it into a normal answer such as \code{reachable: false}.

\subsection{Testing against networkx}

Every tool is tested against networkx on all 300 dataset graphs, and the tools without dataset questions (bipartite,
topological sort) also on random graphs, since only 7 small dataset graphs are bipartite and the dataset has no
directed graphs. The tests check the evidence, not only the yes/no. On 9 October 2026 the whole test suite (439 tests)
passes.

\section{Answers and verifiers}

\code{submit_answer} is built per answer type: \code{number}, \code{yes_no}, \code{node_list}, \code{edge_list}, and
\code{yes_no_with_path} and \code{yes_no_with_cycle}, where a ``yes'' carries the path or cycle in an extra field that
is passed to the verifier. The answer itself remains true or false, so grading does not depend on the evidence.

Each verifier has the signature \code{verify(graph, params, answer, **evidence)} and returns either nothing (accept)
or an error message the model can act on. The verifiers are:
\begin{itemize}
  \item \code{verify_shortest_path}: the path starts and ends at the right nodes, every step is an edge, and its
        length equals the shortest-path distance (see the qualification in Chapter~\ref{ch:design}).
  \item \code{verify_mst}: every edge is a real edge, no edge appears twice, the edges contain no cycle (union-find),
        and there are $n-c$ of them for a graph with $n$ nodes and $c$ components. The dataset's graphs are
        unweighted and some are disconnected, so any spanning forest is a correct answer; the old pipeline's checker
        compared against one exact edge set and therefore marked some correct answers wrong.
  \item \code{verify_connectivity} and \code{verify_cycle}: a ``yes'' must come with a real path or a real simple
        cycle of at least three nodes; a ``no'' is accepted unchecked.
  \item \code{verify_neighbors}: every listed node is a real neighbor and is listed once; a missing neighbor cannot be
        detected.
  \item \code{verify_path_via}: a real route from source to target through the waypoint, whose length equals the sum
        of the two shortest-path distances.
\end{itemize}

\section{Result handles and bounded messages}
\label{sec:handles}

\subsection{Handles}

A value with more than 200 numbers is stored in a per-run \code{HandleStore}, and the model sees a summary with an
identifier, the length and the first five items. A list of large items, such as the components of a graph, gets one
handle per item. Once a handle exists, \code{read_result(handle, offset, limit)} is offered, and \code{submit_answer}
accepts a handle string in place of the value. Before that, the tool list is unchanged, so the current dataset, whose
largest tool output has 98 numbers, never sees handles and earlier results are unaffected.

The first real test illustrates principle 4. On a random graph with 10,000 nodes, 25,000 edges and 68 components,
qwen3:8b was asked for a minimum spanning tree (the answer has 9,932 edges). The tool result came back as a
336-character summary instead of about 110,000 characters, as intended. The model then submitted the handle twice,
was rejected both times, and gave up, explaining that the system required the actual list of edges. The note said
``pass the handle to submit\_answer'', but the schema of \code{submit_answer} said the answer must be an array of
edges, so a string was impossible. After the fix (once a handle exists, the schema becomes ``edge list or a handle
string'', and a handle wrapped in a list is resolved too), the same question took two steps, 35 seconds and 1,974
prompt plus 601 completion tokens, and the verifier checked all 9,932 edges.
% source: docs/learning_log.md (2026-10-07, M3 step 3), results/harness_runs/handles_big_graph_test

\subsection{Nothing the model reads grows with the graph}

Handles only cover tool results. A review found other messages that scaled with the graph: the missing-node error
listed every node identifier, about 59,000 characters on a 10,000-node graph for a single typo. It now reads ``Node
123456 does not exist. G has 10000 nodes, with ids 0 to 9999.'' An audit of every formatted message then produced
Table~\ref{tab:bounded}. The fix is a helper, \code{brief()}, used wherever a value is repeated to the model, and a
test file that checks every such message on a 10,000-node graph.

\begin{table}[ht]
  \centering
  \small
  \caption{Size of messages that repeat a value back to the model, on a 10,000-node graph, before and after the audit.}
  \label{tab:bounded}
  % source: docs/learning_log.md (2026-10-08, missing-node errors and follow-up audit)
  \begin{tabularx}{\textwidth}{>{\raggedright\arraybackslash}X r >{\raggedright\arraybackslash}p{5cm}}
    \toprule
    Message & Before (characters) & After \\
    \midrule
    Wrong-type answer repeated in full & 174,455 & about 280 (first items and size) \\
    A page of large items from \code{read_result} & 69,676 & about 3,500 (each big item its own handle) \\
    The whole cycle repeated by the cycle verifier & 28,923 & names the repeated node \\
    Missing node: every node listed & about 59,000 & 66 \\
    Unknown handle: every stored handle listed & 438 (grows) & first 10 and a count \\
    \bottomrule
  \end{tabularx}
\end{table}

\section{The MCP server}

The MCP server (\code{harness/tools/mcp_server.py}) gives other harnesses exactly our tools. It uses the official
\code{mcp} SDK's low-level server: the tool list is built from the same generated schemas (same names, descriptions
and schemas), calls go through \code{run_tool} (same validation and errors, flagged as errors in MCP), and large results
become handles. The SDK's own argument validation is switched off so that error messages match our loop. The server
speaks stdio, and the client launches one server process per question with the graph file as an argument, so handles
never leak between questions. Differences from our loop: \code{read_result} is listed from the start, and
\code{submit_answer} is not served yet; how a general-purpose harness should submit structured answers is decided in
M6. Our own loop does not use MCP; it calls \code{run_tool} directly.

FastMCP was considered and not used: its main convenience, schemas generated from function signatures, is the opposite
of what is needed, which is our schemas byte for byte. The server is an adapter of about a hundred lines. Nine tests
check that the tool list equals ours, that every tool's result through MCP equals \code{run_tool} on dataset graphs,
that errors match, that handles work, and that the server runs as a real subprocess over stdio.
\todo[inline]{Open decision recorded in CLAUDE.md: keep the mcp package as a project dependency (it is in pyproject.toml now).}

\section{The run\_python sandbox}

\subsection{Design}

\code{run_python} lets the model write code. It is off by default and offered in one of two modes. In
\code{"tools"} mode, the code sees our graph tools as plain Python functions and cannot import networkx, so it
composes our tools instead of bypassing them and their evidence. In \code{"networkx"} mode, the code also gets the
graph \code{G} and \code{nx}: the escape hatch, kept for comparison.

The code runs in a Docker container, one per question, started on the first call (about 0.3 seconds) and reused for
later calls (about 0.3 milliseconds per call), so variables persist between calls. The container has no network, a
read-only file system except a small \code{/tmp}, 512~MB of memory, one CPU, at most 64 processes, no capabilities, a
non-root user, and 10 seconds per call, after which it is killed and restarted and the reply says that variables were
lost. Only the graph file and the \code{harness/} package are mounted, read-only. The model gets back the value of
\code{result} (large values become handles), the printed output (shortened), and errors in the form ``line 2:
IndexError: list index out of range'', pointing at the model's line rather than at harness internals. The backend is
swappable; a WebAssembly backend that needs no Docker is the plan for a public release. On the MacBook, Docker is
provided by OrbStack; on the GPU machine by Docker Engine.
\todo[inline]{Open decision recorded in CLAUDE.md: Docker vs OrbStack/Colima for the sandbox on macOS.}

\subsection{Plain values in code}

The first version exposed the tools inside code exactly as the model sees them in the loop, returning JSON
dictionaries. That failed: \code{for n in get_neighbors(4)} iterated over the dictionary's keys, and an error
dictionary flowed on as if it were a result. Inside code the tools now return plain values, as networkx does
(\code{get_neighbors(4)} returns \code{[2, 5]}; \code{shortest_path(a, b)} returns a path or \code{None}), and tool
errors raise \code{ValueError}. \code{graph_info()} and \code{is_bipartite()} keep their dictionaries, and the tool
description shows what \code{graph_info()} returns. If the model redefines a tool function (one run defined
\code{get_neighbors} in terms of itself and broke every later call), the original is restored after the call, with a
note. The lesson recorded at the time: an interface the model already knows beats a better one it has to learn, so
ours was made to look familiar.

\section{Context window and compaction}

\subsection{The context window}

Ollama serves every model with a 4,096-token context unless asked otherwise, regardless of what the model supports.
Past that limit the prompt is cut and the reply stops short, with no error. This silently damaged one family of
experiments (Section~\ref{sec:case-4k}). The model layer now asks Ollama for 16,384 tokens, logs the context window in
every run's first trace event, and flags any prompt above 90\% of it. Hosted models keep the provider's context.

\subsection{Compaction}

When a prompt passes 75\% of the window, the loop shortens the conversation before the next call, by fixed rules and
without an LLM summary (free and reproducible). Older tool results longer than 300 characters become a short version
with a note (``shortened; call the tool again for the full result''), older \code{run_python} code longer than 200
characters is cut to its start, and the system prompt, the question and the last two rounds stay word for word. No
message is deleted, because the chat format requires every tool call to keep its reply, and handles keep their full
values in the store. Each compaction is logged and counted. If nothing is left to shorten and the prompt is above
90\%, the run ends with status \code{context_full} instead of continuing with a cut prompt. The cost is that the
provider's prompt cache restarts once after a compaction, because the prefix has changed. In the runs reported in
Chapter~\ref{ch:results}, no conversation was compacted.

\section{Prompt versioning}

Every piece of text the harness says to the model lives in \code{harness/prompts.py}, together with
\code{PROMPT_VERSION}, which is logged in every run. The version is bumped, with a changelog line, whenever any
model-facing text changes, including tool descriptions. Table~\ref{tab:prompt-versions} summarises the changelog.
Appendix~\ref{app:prompts} reproduces the current texts.

\begin{table}[ht]
  \centering
  \small
  \caption{Prompt versions (changelog in \code{harness/prompts.py}).}
  \label{tab:prompt-versions}
  % source: harness/prompts.py (PROMPT_VERSION changelog)
  \begin{tabularx}{\textwidth}{l>{\raggedright\arraybackslash}X}
    \toprule
    Version & What changed \\
    \midrule
    v1 & Everything up to and including the M2 pilots \\
    v2 & v1 plus \code{distances_from} and \code{neighborhood}; used by the M2 closing runs and the Claude Code comparison on large graphs \\
    v3 & Merge of the two machines' work: Pydantic schemas, three new tools, result handles, bounded messages, waypoint question \\
    v4 & \code{run_python} description: plain values, errors raise, do not redefine tools \\
    v5 & \code{CODE_HINT} (strategy hint, only with \code{run_python} and the hint on); combine-task wording \\
    v6 & \code{EMPTY_REPLY} nudge; \code{graph_info()} shown in the \code{run_python} description; redefined tools restored \\
    v7 & Compaction texts; runs that never approach the limit see the same text as v6 \\
    \bottomrule
  \end{tabularx}
\end{table}

\section{Logging and analysis}

\subsection{Traces and results}

Every run writes the events \code{run_start} (configuration, prompt version, context window), \code{model_call}
(content, tool calls, thinking, tokens, latency), \code{tool_call} (arguments and the result as the model saw it),
\code{rescued_tool_call}, \code{verifier_rejected}, \code{nudge}, \code{compacted} and \code{run_end} to
\code{traces.jsonl}, and one row per answer to \code{results.jsonl}. An early version of the runner re-read the whole
trace file after every question to find that question's token counts, which made the work grow quadratically with the
number of questions; the loop now returns these counts directly. The log is for later analysis, not for passing data
between parts of the program.

\subsection{The report}

\code{eval/analysis/report.py} puts experiments side by side, overall and per task: lenient and strict accuracy, a 95\%
bootstrap confidence interval over questions (all runs of a question are resampled together), the share of answers a
verifier actually looked at, the number of answers sent back by a verifier, runs without an answer, compactions, and
tokens and time per question. With \code{--code} it reports how \code{run_python} was used, and with \code{--process}
the process metrics below. All tables in Chapter~\ref{ch:results} are produced by this script.

\subsection{Process metrics}
\label{sec:process-metrics}

\code{eval/analysis/process.py} measures whether the model called the right tools with the right arguments, using the
definitions of the GDS Agent benchmark \cite{gdsagentbench}. For each task, \code{EXPECTED_CALLS} lists the calls an
ideal run makes; each expected call is a slot that one of several equivalent tools can fill (for example, a node's
degree from \code{degree} or from \code{get_neighbors}). With $E$ expected calls, $M$ of them made, and $T$ scored tool
calls in the run:
\begin{align*}
  \text{recall} &= M / E, &
  \text{precision} &= M / T, &
  F_1 &= \frac{2PR}{P + R}.
\end{align*}
Repeating a call or exploring lowers precision. Parameter match is the share of right arguments in the expected calls
that were made directly, and exact match means every expected call was made and nothing else. Three deliberate
differences from the original definitions follow from our setup: \code{run_python} counts as one tool call and fills
an expected slot when its code calls that tool function (its arguments cannot be read, so it adds no parameter score);
argument pairs $(u, v)$ and (source, target) also match when swapped, since all graphs are undirected; and the
parameter score takes the best matching call, since the waypoint task makes the same tool call twice. The ending tools
and \code{read_result} are not scored.

\subsection{Trajectory analysis}

For multi-step tasks, \code{eval/analysis/trajectory.py} classifies each run by its graph tool calls and its answer:
ideal, extra calls, alternative path (different calls, correct answer), right calls but wrong answer, or wrong path,
with flags for detours such as taking the shortcut, using \code{has_path} for a leg, or exploring manually with
\code{get_neighbors}. The ``alternative path'' class exists because the first version labelled a correct and well
reasoned run as incomplete; process evaluation has to judge whether each step is justified, not whether it matches one
expected sequence.
